% !PRINTED 27
\weberpage{27}

\begin{center}
\S.\,8.

\gesperrt{Metacyclic Groups.}
\end{center}

We defined metacyclic groups in \S.\,176 of the first volume, in connection with the permutation groups, and there learned their significance for the solution of equations. The concept, however, is independent of the particular meaning of the group elements, and we therefore now define generally:

\gesperrt{A group whose index series consists entirely of primes shall be called a metacyclic group.}

For these groups there holds a theorem important for the applications to algebra, but which is independent of these applications and which we therefore derive here. The theorem reads:

IV.\ \gesperrt{A metacyclic group has a commutative normal
divisor different from the unit
group}\footnote{C. J o r d a n, Traité des substitutions, Livre IV.}.

If $P$ is a metacyclic group and

\weq{1}{
\begin{array}{c}
P,\ P_1,\ P_2 \cdots P_{\mu-1},\ 1\\
j_1,\ j_2 \cdots j_{\mu-1},\ j_\mu
\end{array}
}

is the composition series, whose index series $j_1, j_2 \cdots$ consists entirely of primes, then $P_{\mu-1}$ is in any case a commutative group; for if $\pi$ is any element of $P_{\mu-1}$ different from 1, then the powers

\weq{2}{1,\ \pi,\ \pi^2 \cdots \pi^{j_\mu - 1}}

(because $j_\mu$ is a prime) are all different from one another, and therefore make up the whole group $P_{\mu-1}$; but the cyclic group (2) is certainly commutative, because for any two exponents $h, k$ we always have

\[
\pi^h \pi^k = \pi^k \pi^h = \pi^{h+k}
\]

At the same time, by the definition of the composition series. % [sic] the full stop after "Compositionsreihe" is as printed; a comma is evidently intended
$P_{\mu-1}$ is a normal divisor of $P_{\mu-2}$.

We accordingly now assume that we have a commutative normal divisor $Q_\nu$, different from the unit, of some group $P_\nu$ of the sequence (1), and suppose moreover,
%CONT>>
